\section{Related Work}
\label{sec:related}
Module systems such as Environment Modules and Lmod underpin most HPC software
stacks by organizing compiler, MPI, and library families into hierarchical
module trees.\cite{EnvModules2024,Lmod2011} Package managers like Spack and
EasyBuild automate dependency resolution and software installation across
large software stacks.\cite{Spack2015,EasyBuild2012} These tools are essential
for building and distributing software, but they do not preserve PE-specific
wrapper semantics or the MPI metadata coupling that Cray/HPE systems rely on.

Cray PE documentation emphasizes the role of `ftn/cc/CC` wrappers and PE
metadata paths in maintaining portable build workflows across modules and
targets.\cite{CrayPE2024} Compiler wrapper techniques are widely used to inject
site defaults or architecture flags, yet they typically assume a single
compiler stack rather than a layered environment that overlays experimental
toolchains onto a production PE. Afar-prgenv complements these approaches by
focusing on the user-support problem of early compiler evaluation in a
production PE, with an emphasis on preserving the `ftn/cc/CC` contract and
vendor library integration.

Validation suites such as SOLLVE's OpenMP V\&V provide systematic coverage of
OpenMP semantics and offload behavior,\cite{SOLLVE2020,OpenMP2024} and MPI
implementations commonly ship their own test suites.\cite{MPICH2023} Afar-prgenv
targets a different layer: the interoperability between vendor compiler drops,
PE wrappers, and PE-provided libraries. Its repro suite prioritizes build and
link correctness, while existing validation suites are used after the PE
compatibility gate to assess functional readiness in application workflows.

Finally, site operators often maintain ad hoc module overlays to stage
experimental compilers, but these overlays typically bypass the Cray PE
wrapper contract or require users to hand-curate library flags. Afar-prgenv
formalizes that overlay into a reproducible, metadata-driven workflow that
keeps the wrapper behavior intact while exposing the new toolchain to
existing PE libraries and diagnostics.
